\documentclass{ctexart}
\usepackage{template}
\usepackage{esint,extarrows}
\begin{document}\pagestyle{empty}
\begin{center}\Large
北京大学2025-2026学年第二学期期末试题\\
    高等数学B2
\end{center}
\begin{enumerate}[leftmargin=*,label=\textbf{\arabic*.},topsep=0pt,parsep=0pt,itemsep=0pt,partopsep=0pt]
    \item \textbf{(15\songti{分})}\ 判断无穷级数
\[\sum_{n=3}^{\infty}\frac1{n(\ln n)^p}\qquad(p>0)\]
是否收敛.

    \item \textbf{(15\songti{分})}\ 设
\[a_n=\int_0^{n\pi}x|\sin x|\di x\qquad(n=1,2,\ldots)\]
求级数$\displaystyle\sum_{n=1}^{\infty}a_n/3^n$的和.
（提示：计算$a_n$时可作代换$x=n\pi-t$.）

    \item \textbf{(15\songti{分})}\ 解答下列问题.
\begin{enumerate}[label=\tbf{(\arabic*)},topsep=0pt,parsep=0pt,itemsep=0pt,partopsep=0pt]
\item 求$f(x)=\arctan\dfrac{2-2x}{1+4x}$在$x=0$处的Taylor级数.（提示：先求导数.）
\item 求幂级数$\displaystyle\sum_{n=0}^{\infty}(-1)^n\dfrac{x^{2n}}{(2n)!!}$的和函数,其中$0!!=1$.
\end{enumerate}

    \item \textbf{(15\songti{分})}\ 判断瑕积分
\[\int_0^1\frac{\sin(1/x)}{x^p}\di x\qquad(0<p\leqslant2)\]
的敛散性;若收敛,判断其为条件收敛还是绝对收敛.

    \item \textbf{(15\songti{分})}\ 将$f(x)=2+|x|$ $(-1\leqslant x\leqslant1)$作周期为$2$的Fourier展开,并用此展开式求$\displaystyle\sum_{n=1}^{\infty}1/(2n-1)^2$的和.

    \item \textbf{(15\songti{分})}\ 计算积分
\[I=\int_0^{+\infty}\frac{\e^{-\alpha x}-\e^{-\beta x}}x\sin x\di x,
\qquad 0<\alpha<\beta\]

    \item \textbf{(10\songti{分})}\ 判断级数
\[\sum_{n=1}^{\infty}\left(\frac1n-\ln\frac{n+1}{n}\right)\]
的敛散性,再利用所得结论判断极限$\displaystyle\lim_{n\to\infty}\dfrac{1+1/2+\cdots+1/n}{\ln n}$是否存在,并求其值.

\end{enumerate}
\end{document}
